\documentclass[11pt]{article}
\usepackage[margin=2.3cm]{geometry}
\usepackage{amsmath,amssymb,amsthm,booktabs}
\newtheorem{theorem}{Theorem}
\newtheorem{lemma}[theorem]{Lemma}
\newtheorem{corollary}[theorem]{Corollary}
\theoremstyle{remark}
\newtheorem{remark}[theorem]{Remark}
\newcommand{\hG}{h_\Gamma}
\newcommand{\VO}{V_O}
\newcommand{\Pih}{\Pi_h}
\newcommand{\Lk}{\mathcal L}
\newcommand{\dn}{\partial_\nu}
\newcommand{\norm}[1]{\|#1\|}
\newcommand{\Ql}{Q_{\rm lock}}
\title{Strong no-slip for Scott--Vogelius with locked wall vertices:\\ closing S4 (notes/v6/strong2)}
\date{2026-09-28}
\begin{document}
\maketitle

\section*{Summary and status}
Notation as in \texttt{paper/sections/02\_setting.tex} and \texttt{notes/v6/strong/REPORT.md}.
A \emph{locked} vertex is a vertex $z$ of $\Gamma_h$ lying in exactly two triangles; $\Lk$ is the set of
locked vertices. $\phi_z$ is the turning angle of $P_h$ at $z$; for the inscribed polygon
$\phi_z=\arcsin(|e_1|/2)+\arcsin(|e_2|/2)$, so $c_0\hG\le\phi_z\le 1.1\,\hG$ by (M0).

\begin{center}\small\begin{tabular}{p{3.2cm}p{9.6cm}p{2.4cm}}
\toprule
Claim & Content & Status\\\midrule
Lemma~\ref{lem:K} & vertex gradients of $\VO$ at a locked $z$: exactly $\{(c\otimes g_1,c\otimes g_2)\}$; exact inverse formulas & PROVED\\
Theorem~\ref{thm:L} & uniform right inverse of $\div$ on $\Ql=\{q:\ q|_{T_1}(z)=q|_{T_2}(z)\ \forall z\in\Lk\}$ & PROVED$^*$\\
Corollary~\ref{cor:L} & hypothesis (L) of the S4 note (constraint $q|_{T_i}(z)=0$) & PROVED$^*$\\
Theorem~\ref{thm:sharp} & $\min\{\norm{\nabla v}:\div v=q\}\asymp\norm q+\big(\sum_z(h_z|q_1(z)-q_2(z)|/\sin\phi_z)^2\big)^{1/2}$; $\beta_h\asymp\min_z\phi_z\asymp\hG$ & PROVED$^*$\\
Theorem~\ref{thm:S4} & $\norm{\nabla(\tilde u-u_h)}\le C(h^k+\hG^{3/2}+\hG(\sum_{\Lk}|\dn u(z)|^2)^{1/2})$ & PROVED$^*$\\
Corollary~\ref{cor:two} & two-sided rate (with S2, S3$'$) & PROVED$^*$\\
Theorem~\ref{thm:P}(i),(ii),(iii) & pressure: $\Ql$-part at velocity rate; raw error $\le CR/\phi_{\min}$; filtered pressure & PROVED$^*$\\
Sharpness of Thm~\ref{thm:P}(ii) & raw pressure error $O(1)$ (one lock), $\sim\hG^{-1/2}$ (mesh (a)) & NUMERICAL\\
Adjacency & no separation of locked vertices needed & PROVED$^*$ + NUMERICAL\\
\bottomrule\end{tabular}\end{center}
$^*$``PROVED'' modulo the following cited standard ingredients. None of them is applied on a locked star in a
way that could depend on $\phi$:
(a) the uniform continuous inf-sup constant on $\Omega_h$ (proof of \texttt{lem:infsup} in the paper, Bogovski\u{\i}/Galdi);
(b) the $P_2$ Fortin operator with $\int_e\Pi_F w=\int_e w$ on edges, stable under shape regularity
(Crouzeix--Raviart 1973; Boffi--Brezzi--Fortin 2013, \S8.4);
(c) the element bubble lemma (Lemma~\ref{lem:bubble}), classical (Scott--Vogelius 1985) and checked exactly here for $k=3,\dots,8$.
The vertex lemma (Lemma~\ref{lem:V}) is proved below (surjectivity by a dimension count, uniformity by compactness),
so we do not rely on the internal structure of the Guzm\'an--Scott proof.

\paragraph{Answer to the question ``which constraints are necessary''.}
At a locked vertex there is \emph{no} exact constraint on $\div\VO$: the vertex-divergence map
$\VO\to\mathbb R^2$, $v\mapsto(\div v|_{T_1}(z),\div v|_{T_2}(z))$ is onto (the vertex is nearly singular,
not singular), so $\div\VO=\Pih\cap L^2_0$. What fails is \emph{uniformity}: the preimage of
$(q_1,q_2)$ costs $\asymp h_z|q_1-q_2|/\sin\phi_z+h_z|q|$ in $H^1$ (Theorem~\ref{thm:sharp}). The single
scalar condition $q|_{T_1}(z)=q|_{T_2}(z)$ -- the $\phi\to0$ limit of the Scott--Vogelius alternating-sum
condition, which is exact at a flat 2-triangle boundary vertex -- is necessary and sufficient for a
$\phi$-independent bound. The two conditions $q|_{T_1}(z)=q|_{T_2}(z)=0$ of hypothesis (L) are more
than needed; they are what the approximation step naturally produces (the div-free solution has $\nabla u_h(z)=0$).

\section{Hypotheses}
\textbf{(H)} (M0), (M1) of the paper; $k\ge4$; every vertex of $\Gamma_h$ has at least two triangles
(implied by (M1)); (M2) with constant $\Theta_0$ holds at every vertex \emph{not} in $\Lk$
(interior vertices, vertices of $\partial R$, and $\Gamma_h$ vertices with $m\ge3$ -- for the latter it is
automatic by L2 of the S4 note when $\phi_z\le\theta_{\min}/2$).
No condition on the number or arrangement of locked vertices. In particular locked vertices may be
adjacent (two locked neighbours $z,z_+$ share the triangle $[z,z_+,w]$ and the apex $w$), and the apex $w$ of
several locked stars is an ordinary vertex subject to (M2).
Constants $C,c$ depend only on $k$, $\theta_{\min}$, $c_0$, $\Theta_0$, $L$ (and on norms of $u,p$ where
displayed); never on $h,\hG,\phi_z,\#\Lk$ or the positions of locked vertices.

\section{Local structure at a locked vertex}
Let $z\in\Lk$ with triangles $T_1=[z_-,z,w]$, $T_2=[z,z_+,w]$, wall edges $e_1=[z_-,z]$, $e_2=[z,z_+]$ and
interior edge $s=[z,w]$. Let $\lambda_x$ be the global hat function of vertex $x$ and put
$g_i:=\nabla\lambda_w|_{T_i}$. Then $g_i=\nu_i/H_i$ with $\nu_i$ the unit normal of $e_i$ pointing into
$T_i$ and $H_i=\mathrm{dist}(w,\mathrm{line}\,e_i)\asymp h_z:=\mathrm{diam}(T_1\cup T_2)\asymp\hG$.
Since $\lambda_w$ is linear on $s$ from $0$ to $1$,
\begin{equation}\label{eq:gs}
g_1\cdot(w-z)=g_2\cdot(w-z)=1 .
\end{equation}
Let $D_z$ be the $2\times2$ matrix with rows $g_1^{\mathsf T},g_2^{\mathsf T}$.

\begin{lemma}[Locked star]\label{lem:K}
\begin{enumerate}
\item[(a)] For every $v\in\VO$ there is a unique $c_v\in\mathbb R^2$ with $\nabla v|_{T_i}(z)=c_v\otimes g_i$
($i=1,2$). Hence $(\div v|_{T_1}(z),\div v|_{T_2}(z))=D_zc_v$, and $|c_v|\le C\norm{\nabla v}_{T_1}$.
\item[(b)] For $c\in\mathbb R^2$, $\delta\in\mathbb R$ put
$K_z(c,\delta):=c\,\lambda_w\lambda_z^2+\delta\,n_s\lambda_z^2\lambda_w^2$ ($n_s$ the unit normal of $s$ from $T_1$ to $T_2$).
Then $K_z(c,\delta)\in\VO$, $\operatorname{supp}K_z\subset T_1\cup T_2$, $\nabla K_z=0$ at $z_-,z_+,w$ (indeed at every vertex other than $z$),
$\nabla K_z|_{T_i}(z)=c\otimes g_i$, $\norm{\nabla K_z}\le C(|c|+|\delta|)$, and
$\int_{T_1}\div K_z=-\int_{T_2}\div K_z=|s|\big(\tfrac1{12}c\cdot n_s+\tfrac1{30}\delta\big)$.
\item[(c)] $|\det D_z|=\sin\phi_z/(H_1H_2)$; in particular $D_z$ is invertible.
\item[(d)] $D_z(w-z)=(1,1)^{\mathsf T}$. Hence the vertex data $(\sigma,\sigma)$ are realised by
$K_z(\sigma(w-z),0)=\sigma(w-z)\lambda_w\lambda_z^2$, which has $\norm{\nabla\cdot}\le C|\sigma|h_z$ and
\emph{zero} flux through $s$ (so zero mean divergence on $T_1$ and on $T_2$).
\item[(e)] $|D_z^{-1}(1,0)^{\mathsf T}|=H_1/\sin\phi_z$ and $|g_1-g_2|\le C\phi_z/h_z$.
\end{enumerate}
\end{lemma}
\begin{proof}
(a) $v=0$ on $e_1$, so $\nabla v|_{T_1}(z)t_1=0$ and $\nabla v|_{T_1}(z)=a\otimes\nu_1=c\otimes g_1$; likewise
$\nabla v|_{T_2}(z)=c'\otimes g_2$. Continuity across $s$ gives $\nabla v|_{T_1}(z)(w-z)=\nabla v|_{T_2}(z)(w-z)$,
i.e.\ $c=c'$ by \eqref{eq:gs}. The bound: $|c_v||g_1|=|\nabla v|_{T_1}(z)|\le Ch_z^{-1}\norm{\nabla v}_{T_1}$ (inverse
inequality) and $|g_1|=1/H_1\ge c/h_z$.
(b) $\lambda_w\lambda_z$ is nonzero only on triangles containing both $z$ and $w$, i.e.\ $T_1\cup T_2$; it is
continuous (product of hats) of degree $3$ resp.\ $4\le k$. $\lambda_w=0$ on $e_1,e_2$ and $\lambda_z=0$ on
$[z_-,w],[w,z_+]$, so $K_z\in\VO$. At a vertex $x\ne z$, $\lambda_z(x)=0$ and $\nabla(\lambda_z^2)(x)=0$, so
$\nabla K_z(x)=0$. At $z$, $\lambda_z=1,\lambda_w=0$ give $\nabla K_z|_{T_i}(z)=c\otimes g_i$ (the $\delta$-term has a
factor $\lambda_w^2$). Fluxes: $K_z=0$ on $\partial(T_1\cup T_2)$, and on $s$ (param.\ $\lambda_w=r$,
$\lambda_z=1-r$) $\int_s\lambda_w\lambda_z^2=|s|/12$, $\int_s\lambda_z^2\lambda_w^2=|s|/30$. The $H^1$ bound is scaling.
(c) $g_1\times g_2=\sin\angle(\nu_1,\nu_2)/(H_1H_2)$ and the normals of consecutive chords make the angle $\phi_z$.
(d) is \eqref{eq:gs}; the flux is $\sigma(w-z)\cdot n_s\,|s|/12=0$ since $w-z\parallel s$.
(e) $D_z^{-1}(1,0)^{\mathsf T}=(\det D_z)^{-1}g_2^{\perp}$, of length $|g_2|/|\det D_z|=H_1/\sin\phi_z$. For the second
claim let $\alpha_i$ be the angle of $T_i$ at $z$; $\alpha_1+\alpha_2=\pi+\phi_z$ and $H_i=|s|\sin\alpha_i$, so
$|H_1-H_2|=|s||\sin\alpha_1-\sin(\alpha_1-\phi_z)|\le|s|\phi_z$, and
$|g_1-g_2|\le|\nu_1-\nu_2|/H_1+|H_1^{-1}-H_2^{-1}|\le C\phi_z/h_z$.
\end{proof}
Exact rational checks of (b)--(e) on inscribed stars with chord parameter $t=1/10,\dots,1/160$ (two symmetric apex
positions) and $t\le1/40$ (an off-axis apex): \texttt{local\_L.py}, log \texttt{local\_L.log}. With $h=|e_1|$:
$h\sigma_{\min}(D_z)/\phi_z\to0.943$ resp.\ $0.471$; $|c|/h$ stays bounded ($0.75$, $1.50$, $0.82$) for data $(1,1)$;
$|c|\phi/h$ stays bounded ($0.75$, $1.50$, $0.74$) for data $(1,0)$, so there $|c|\sim h/\phi$. The solution for $(1,1)$ is
\emph{exactly} $c=w-z$ (asserted in the script). Singular values use floating point on the exact matrix.

\section{The constrained inf-sup condition}
\begin{lemma}[Element bubbles]\label{lem:bubble}
For every triangle $T$ and $q\in P_{k-1}(T)$ with $\int_Tq=0$ and $q=0$ at the three vertices there is
$v\in P_k(T)^2\cap H^1_0(T)^2$ with $\div v=q$ and $\norm{\nabla v}_T\le C(\theta_{\min},k)\norm q_T$.
\end{lemma}
\begin{proof}
On the reference triangle $\div$ maps the bubbles $(b_TP_{k-3})^2$ into the target space (a divergence of an
$H^1_0$ field has zero mean, and $\nabla v=0$ at vertices because $v$ vanishes on two independent edges).
Surjectivity for $k=3,\dots,8$ is an exact rank computation (\texttt{local\_L.py} (B): ranks $2,6,11,17,24,32$
equal the target dimensions); for general $k$ it is classical (Scott--Vogelius 1985). Transport by the
contravariant Piola map, which preserves $H^1_0\cap P_k$ and scales $\div$ by $1/\det B_T$; the constant depends
on shape regularity only.
\end{proof}

\begin{lemma}[Vertex lemma]\label{lem:V}
Let $y\notin\Lk$ be a vertex with triangles $T_1,\dots,T_m$ and $\Theta(y)\ge\Theta_0$. For every
$r\in\mathbb R^m$ there is $v_y\in\VO$ with $\operatorname{supp}v_y\subset\mathrm{star}(y)$ such that
(i) $\div v_y|_{T_j}(y)=r_j$; (ii) $\nabla v_y(x)=0$ at every vertex $x\ne y$; (iii) $\int_T\div v_y=0$ for every $T$;
(iv) $\norm{\nabla v_y}\le Ch_y|r|$.
\end{lemma}
\begin{proof}
Let $x_j$ be the neighbours of $y$. Set $\ell=\sum_jc_j\lambda_{x_j}$ with $c_j\in\mathbb R^2$ and $c_j=0$ when
$[y,x_j]\subset\partial\Omega_h$, and $v_y=\lambda_y^2\ell+\sum_e\delta_e n_e\lambda_y^2\lambda_{x_e}^2$, the sum over the
edges $e=[y,x_e]$ of the star not on $\partial\Omega_h$. Every term has the factor $\lambda_y^2$, so $v_y$ is supported in the
star, vanishes on its outer boundary, and $\nabla v_y(x)=0$ at all vertices $x\ne y$: (ii). On a boundary edge
$[y,x_0]$ of the star, $\ell=c_0\lambda_{x_0}=0$ and $\lambda_{x_e}=0$ for $x_e\ne x_0$, so $v_y\in\VO$.
At $y$, $\nabla v_y|_{T}(y)=\nabla\ell|_T$ (the bubbles have $\lambda_{x_e}^2$).

\emph{Surjectivity.} The tuples $(\nabla\ell|_{T_j})_j$ are exactly the vertex-gradient tuples of $C^0$ fields
vanishing at $y$ (and on the boundary edges through $y$): dimension $2m$ (interior) resp.\ $2(m-1)$ (boundary).
The kernel of the trace map on this set is the set $A_y$ of \emph{div-free} tuples; by the argument of L1 in the S4
note (write $G=JH$, $H$ symmetric, jumps $\beta_is_i^\perp s_i^{\perp\mathsf T}$, closure relation of rank
$\min(3,L_y)$, $L_y$ = number of distinct edge lines), $\dim A_y=m+3-\min(3,L_y)$ (interior) resp.
$m+1-\min(3,L_y)$ (boundary). Hence the rank of $c\mapsto(\div\ell|_{T_j})_j$ is $m-3+\min(3,L_y)$ in both
cases, i.e.\ it is onto $\mathbb R^m$ iff $L_y\ge3$. If $L_y\le2$ then consecutive angles are supplementary and
$\Theta(y)=0$; so $\Theta(y)\ge\Theta_0$ implies surjectivity.

\emph{Uniformity.} Normalise $h_y=1$. The set of stars with $m\le2\pi/\theta_{\min}$, all angles
$\ge\theta_{\min}$, edge ratios $\le C(\theta_{\min})$ and $\Theta\ge\Theta_0$ is compact, and the least singular value
of the (onto) map depends continuously on the star, so it has a positive minimum. Rescaling gives $|c|\le Ch_y|r|$.

\emph{Means.} $\sum_j\int_{T_j}\div(\lambda_y^2\ell)=\int_{\partial\,\mathrm{star}}\lambda_y^2\ell\cdot n=0$. The edge bubble on
$e$ moves flux $\delta_e|e|/30$ between the two triangles sharing $e$; the triangles of the star with these edges
form a cycle (interior $y$) or a path (boundary $y$), which is connected, so every zero-sum vector of triangle
means is matched, with $|\delta_e|\le C\sum_j|\int_{T_j}\div(\lambda_y^2\ell)|/h_y\le C|c|$. This gives (iii), (iv).
\end{proof}

\begin{theorem}[Uniform constrained inf-sup]\label{thm:L}
Assume (H). Let
\[\Ql:=\{q\in\Pih\cap L^2_0(\Omega_h):\ q|_{T_1(z)}(z)=q|_{T_2(z)}(z)\ \text{for all }z\in\Lk\}.\]
For every $q\in\Ql$ there is $y\in\VO$ with $\div y=q$ and $\norm{\nabla y}\le C\norm q$, $C$ independent of
$h,\hG,(\phi_z)_{z\in\Lk}$ and of the number and arrangement of locked vertices.
\end{theorem}
\begin{proof}
\emph{Step A (means; vertex gradients at locked vertices removed).}
By (a) there is $w\in H^1_0(\Omega_h)^2$ with $\div w=q$, $\norm{\nabla w}\le C\norm q$. Let $v_A=\Pi_Fw\in\VO$ be its
$P_2$ Fortin interpolant (b): $\int_T\div v_A=\int_Tq$ for all $T$, $\norm{\nabla v_A}\le C\norm{\nabla w}$. Put
\[v_A'=v_A-\sum_{z\in\Lk}K_z\big(c_{v_A,z},-\tfrac52\,c_{v_A,z}\cdot n_{s_z}\big).\]
By Lemma~\ref{lem:K}(b) each subtracted field is flux-neutral on every triangle, has zero gradient at every vertex
other than its own $z$, and kills the vertex gradient of $v_A$ at $z$. Hence $\nabla v_A'(z)=0$ for $z\in\Lk$,
$\int_T\div v_A'=\int_Tq$, and since a triangle lies in at most two locked stars,
$\norm{\nabla v_A'}\le C\norm{\nabla v_A}$ (Lemma~\ref{lem:K}(a)). Let $q_1=q-\div v_A'$; then
$\int_Tq_1=0$ for all $T$ and $q_1|_{T_i(z)}(z)=q|_{T_i(z)}(z)=:\sigma_z$ ($i=1,2$) for $z\in\Lk$.

\emph{Step B (locked vertices).} $v_B:=\sum_{z\in\Lk}\sigma_z(w_z-z)\lambda_{w_z}\lambda_z^2$. By
Lemma~\ref{lem:K}(d) it matches the vertex divergences $(\sigma_z,\sigma_z)$, has zero triangle means, zero gradient
at all other vertices, and $\norm{\nabla v_B}^2\le C\sum_z\sigma_z^2h_z^2\le C\norm{q_1}^2$ (inverse inequality
$h_T|p(x)|\le C\norm p_T$). $q_2:=q_1-\div v_B$ has zero triangle means and vanishes at locked vertices.

\emph{Step C (other vertices).} $v_C:=\sum_{y\notin\Lk}v_y$ from Lemma~\ref{lem:V} with $r_T=q_2|_T(y)$. By (ii) the
vertex values at other vertices and by (iii) the means are unchanged, and
$\norm{\nabla v_C}^2\le C\sum_y h_y^2|r_y|^2\le C\norm{q_2}^2$. $q_3:=q_2-\div v_C$ has zero mean on every
triangle and vanishes at every vertex of every triangle.

\emph{Step D (elements).} Lemma~\ref{lem:bubble} on each $T$ gives $v_D$ with $\div v_D=q_3$,
$\norm{\nabla v_D}\le C\norm{q_3}$.

$y=v_A'+v_B+v_C+v_D$ satisfies $\div y=q$, $y\in\VO$, $\norm{\nabla y}\le C\norm q$. The only place a locked vertex
enters is Lemma~\ref{lem:K}(a),(b),(d), whose constants do not involve $\phi_z$.
\end{proof}

\begin{corollary}[Hypothesis (L) of the S4 note]\label{cor:L}
(L) holds: for $q\in\Pih\cap L^2_0$ with $q|_{T_1(z)}(z)=q|_{T_2(z)}(z)=0$ at every $z\in\Lk$ there is $y\in\VO$,
$\div y=q$, $\norm{\nabla y}\le C\norm q$.
\end{corollary}

\begin{theorem}[Sharp stability of $\div$ without the constraint]\label{thm:sharp}
Assume (H). For $q\in\Pih\cap L^2_0$ write $\ell_z(q):=q|_{T_1(z)}(z)-q|_{T_2(z)}(z)$. Then
\[
c\,N(q)\le\min\{\norm{\nabla v}:v\in\VO,\ \div v=q\}\le C\,N(q),\qquad
N(q):=\norm q+\Big(\sum_{z\in\Lk}\Big(\frac{h_z|\ell_z(q)|}{\sin\phi_z}\Big)^2\Big)^{1/2}.
\]
Consequently the inf-sup constant $\beta_h$ of $(\VO,\Pih\cap L^2_0)$ satisfies
$c\min(1,\min_{\Lk}\phi_z)\le\beta_h\le C\min_{\Lk}\phi_z$, i.e.\ $\beta_h\asymp\hG$ as soon as $\Lk\ne\emptyset$,
and $\div\VO=\Pih\cap L^2_0$ always.
\end{theorem}
\begin{proof}
\emph{Upper.} Repeat the proof of Theorem~\ref{thm:L}, replacing Step B by $v_B=\sum_zK_z(c_z,\delta_z)$ with
$c_z=D_z^{-1}(q_1|_{T_1}(z),q_1|_{T_2}(z))^{\mathsf T}$ and $\delta_z=-\frac52c_z\cdot n_s$. Split the data as
$\sigma(1,1)+\ell_z(q)(1,0)$ with $\sigma=q_1|_{T_2}(z)$; by Lemma~\ref{lem:K}(d),(e)
$|c_z|\le C h_z|\sigma|+H_1|\ell_z(q)|/\sin\phi_z$ (note $\ell_z(q_1)=\ell_z(q)$).
\emph{Lower.} $\norm{\nabla v}\ge\norm{\div v}/\sqrt2=\norm q/\sqrt2$. By Lemma~\ref{lem:K}(a),(e),
$c_v=\sigma(w-z)+\ell_z(q)D_z^{-1}(1,0)^{\mathsf T}$ with $\sigma=q|_{T_2}(z)$, so
$\norm{\nabla v}_{T_1(z)}\ge c|c_v|\ge c\,h_z|\ell_z(q)|/\sin\phi_z-C\norm q_{T_2(z)}$. Square, sum (each triangle in at
most two locked stars) and combine with the first bound.
\emph{Inf-sup.} The lower bound on $\beta_h$ follows from $N(q)\le C\norm q/\min(1,\sin\phi_{\min})$
($h_z|\ell_z(q)|\le C\norm q_{T_1\cup T_2}$). For the upper bound let $\rho_z:=R_1-R_2$, where $R_i\in P_{k-1}(T_i)$ is the
$L^2(T_i)$ Riesz representer of $p\mapsto p(z)$. Then $\int\rho_z=1-1=0$, $(q,\rho_z)=\ell_z(q)$ for $q\in\Pih$,
$\norm{\rho_z}^2=R_1(z)+R_2(z)\ge c/h_z^2$, and for $v\in\VO$, $(\rho_z,\div v)=\ell_z(\div v)=c_v\cdot(g_1-g_2)$,
so $|(\rho_z,\div v)|\le C\phi_zh_z^{-1}\norm{\nabla v}$ by Lemma~\ref{lem:K}(a),(e). Hence
$\sup_v(\rho_z,\div v)/(\norm{\rho_z}\norm{\nabla v})\le C\phi_z$.
\end{proof}

\section{The velocity estimate}
Method (S): $u_h\in V_h$, $u_h|_{\partial R}=g_I$, $u_h|_{\Gamma_h}=0$, $p_h\in\Pih\cap L^2_0$ with
$a_h(u_h,v)-(p_h,\div v)=(\tilde f,v)$ for $v\in\VO$ and $\div u_h=0$; outer flux $m_R=0$.
By Theorem~\ref{thm:sharp}, $(u_h,p_h)$ exists and is unique under (H) (for every $h$, but with a pressure
stability constant $\asymp\hG^{-1}$ -- Section~\ref{sec:p}).

\begin{theorem}[S4]\label{thm:S4}
Assume (H), (D), $m_R=0$. Then
\[\norm{\nabla(\tilde u-u_h)}\le C\Big(h^k+\hG^{3/2}+\hG\Big(\sum_{z\in\Lk}|\dn u(z)|^2\Big)^{1/2}\Big).\]
\end{theorem}
\begin{proof}
\emph{Interpolant.} $w_h:=I_h\tilde u-\sum_{x\in\mathcal N_\Gamma}\tilde u(x)\varphi_x$ as in S1:
$\norm{\nabla(\tilde u-w_h)}\le C(h^k+\hG^{3/2})$, $w_h=0$ on $\Gamma_h$, $w_h=g_I$ on $\partial R$. At $z\in\Lk$,
$T_1,T_2$ carry a $\Gamma_h$ edge, so $h_{T_i}\asymp\hG$ by (M0) and shape regularity. Since $z\in\Gamma$,
$\nabla\tilde u(z)=\nabla u(z)=\dn u(z)\otimes\nu$ (\texttt{lem:wall}); the interpolation error at $z$ is
$\le Ch_T^k$, and zeroing costs $\sum_x|\tilde u(x)||\nabla\varphi_x(z)|\le C\hG^2\cdot\hG^{-1}$. Hence
$|\nabla w_h|_{T_i}(z)|\le|\dn u(z)|+C\hG$ and, by Lemma~\ref{lem:K}(a), $|c_{w_h,z}|\le C\hG(|\dn u(z)|+C\hG)$.
Put $w_h':=w_h-\sum_{z\in\Lk}K_z(c_{w_h,z},0)$. Then $w_h'$ has the same traces, $\nabla w_h'(z)=0$ on both
triangles at every locked $z$, and (supports overlap at most twice, $\#\Lk\le C/\hG$)
\[\norm{\nabla(w_h'-w_h)}^2\le C\sum_{z\in\Lk}\hG^2(|\dn u(z)|^2+\hG^2)\le C\hG^2\sum_{\Lk}|\dn u(z)|^2+C\hG^3.\]
\emph{Divergence.} $q:=\div w_h'\in\Pih$ has mean $m_R=0$, vanishes at every locked vertex on both triangles,
so $q\in\Ql$; and $\norm q\le\sqrt2\norm{\nabla(w_h'-\tilde u)}$ since $\div\tilde u=0$.
Theorem~\ref{thm:L} gives $y\in\VO$ with $\div y=q$, $\norm{\nabla y}\le C\norm q$, so $z_h:=w_h'-y$ is
admissible and divergence-free with $\norm{\nabla(\tilde u-z_h)}\le C\norm{\nabla(\tilde u-w_h')}$.
\emph{Energy argument.} Exactly S1 steps 4--5: $\varepsilon=u_h-z_h\in Z_O$ is pointwise divergence-free, so the
pressure drops out, $|(F,\varepsilon)|\le C\hG^2\norm{\nabla\varepsilon}$ (\texttt{lem:ext}), and Korn gives
$\norm{\nabla\varepsilon}\le C(\norm{\nabla(\tilde u-z_h)}+\hG^2)$.
\end{proof}

\begin{corollary}[Two-sided rate]\label{cor:two}
Under (H), (D), $m_R=0$, $h^k\le\hG^{3/2}$ and $\norm{\dn u}_{L^2(\Gamma)}>0$:
\[\norm{\nabla(\tilde u-u_h)}\asymp\hG^{3/2}+\hG\Big(\sum_{z\in\Lk}|\dn u(z)|^2\Big)^{1/2}\]
for $\hG$ small, by Theorem~\ref{thm:S4}, S2 ($\ge c\hG^{3/2}$) and S3$'$ ($\ge c\hG(\sum_\Lk|\dn u|^2)^{1/2}-C\hG^{3/2}$).
In particular: one lock $\Rightarrow$ rate $1$; a positive arc fraction of locks where $\dn u\ne0$ $\Rightarrow$ rate $1/2$;
$n$ locks $\Rightarrow$ $\hG^{3/2}+\sqrt n\,\hG$.
\end{corollary}

\section{Pressure}\label{sec:p}
Let $\pi_h:=p_h-\Pi(\tilde p-\bar p)\in\Pih\cap L^2_0$ ($\Pi$ = $L^2$ projection) and
$R:=\norm{\nabla(\tilde u-u_h)}+C\hG^2$. As in S1 step 6, $|(\pi_h,\div v)|\le CR\norm{\nabla v}$ for all $v\in\VO$.
Let $P_{\rm lock}$ be the $L^2$-orthogonal projection of $\Pih\cap L^2_0$ onto $\Ql$; its complement is
$\Ql^\perp=\mathrm{span}\{\rho_z:z\in\Lk\}$ (the $\rho_z$ of Theorem~\ref{thm:sharp}, local on the locked stars).

\begin{theorem}[Pressure]\label{thm:P}
Assume (H), (D), $m_R=0$.
\begin{enumerate}
\item[(i)] $\norm{P_{\rm lock}\pi_h}\le CR$.
\item[(ii)] $\norm{\pi_h}\le CR/\min(1,\phi_{\min})\le CR/\hG$.
\item[(iii)] The \emph{filtered pressure} $p_h^f:=P_{\rm lock}p_h$ (computable: subtract $\sum_z\gamma_z\rho_z$, a banded
Gram solve) satisfies $\norm{p_h^f-\Pi(\tilde p-\bar p)}\le C(R+\hG^{k+1/2})$.
\end{enumerate}
\end{theorem}
\begin{proof}
(i) Theorem~\ref{thm:L} gives $v$ with $\div v=P_{\rm lock}\pi_h$, $\norm{\nabla v}\le C\norm{P_{\rm lock}\pi_h}$, and
$(\pi_h,\div v)=\norm{P_{\rm lock}\pi_h}^2$. (ii) Same with Theorem~\ref{thm:sharp}. (iii)
$p_h^f-\Pi(\tilde p-\bar p)=P_{\rm lock}\pi_h-(I-P_{\rm lock})\Pi\tilde p$ (constants lie in $\Ql$). For any $g^*\in\Ql$,
$\norm{(I-P_{\rm lock})\Pi\tilde p}\le\norm{\Pi\tilde p-g^*}$. Take $g^*=\Pi\tilde p-\sum_z\ell_z(\Pi\tilde p)e_z$ with
$e_z\in P_{k-1}(T_1(z))$, $e_z(z)=1$, $e_z=0$ at the other two vertices, $\int e_z=0$, $\norm{e_z}\le Ch_z$ (exists for
$k-1\ge2$); then $g^*\in\Ql$ even for adjacent locks, and $|\ell_z(\Pi\tilde p)|=|(\Pi\tilde p-\tilde p)|_{T_1}(z)-(\Pi\tilde p-\tilde p)|_{T_2}(z)|\le Ch_z^k$.
Sum over $\#\Lk\le C/\hG$.
\end{proof}

\begin{remark}[Sharpness of (ii) -- NUMERICAL]
(ii) is attained: with one lock the raw pressure error does \emph{not} converge ($\norm{\pi_h}\to\approx12$), with
mesh (a) it grows like $\hG^{-1/2}$, and $\norm{\pi_h}_{L^\infty(\text{layer})}\propto\hG^{-1}$ in both cases; the error
is almost entirely the $\Ql^\perp$ component (Table~\ref{tab:filter}, column $\norm{(I-P_{\rm lock})p_h}$). The
filtered pressure converges at the velocity rate (rate $\approx1.1$ with one lock, $\approx0.57$ on mesh (a)).
Heuristically: the shear mode $\dn u_t$ that the lock suppresses is nearly div-free (divergence $O(\phi_z)$), and the
momentum residual it leaves, of size $\asymp\hG|\dn u(z)|$ per lock, can only be balanced by the pressure mode
$\rho_z$ at amplitude $\asymp\hG|\dn u(z)|/\phi_z=O(1)$. A rigorous lower bound for the raw pressure error was not
attempted.
\end{remark}

\section{Adjacent locked vertices}
Nothing in Sections 3--6 uses separation: the correction fields at different locked vertices have zero gradient
at each other's vertices (Lemma~\ref{lem:K}(b)), are flux-neutral per triangle, and each triangle lies in at most
two locked stars (a triangle has at most two $\Gamma_h$ vertices). The only requirement is (M2) at the non-locked
vertices, in particular at the common apex $w$ of a run of consecutive locked vertices; with minimum angle
$\theta_{\min}$ such a run has length $\le 2\pi/\theta_{\min}$ automatically. No hypothesis ``locked vertices are
separated'' is needed. Checked numerically on mesh (a) (alternating) and on a mesh \texttt{fan} with two
\emph{consecutive} locked vertices sharing an apex (Table~\ref{tab:infsup}).

\section{Numerics}
Scripts in \texttt{notes/v6/strong2/}; code base \texttt{code/svn.py}, \texttt{code/strong\_bc.py} ($k=4$).
\begin{itemize}
\item \texttt{infsup.py MESH N [--lobpcg]}: $\beta^2=\min(BA^{-1}B^{\mathsf T}q,q)/(Mq,q)$ on $\VO$ ($A$ = vector
Laplacian) over $\Pih\cap L^2_0$ (\texttt{full}), $\Ql$ (\texttt{lockw}) and the (L)-space (\texttt{locks}); dense
generalized eigensolver for $N\le32$, LOBPCG with $M$-orthogonality constraints otherwise (cross-checked at
$N=16,24$ on mesh \texttt{one}). Logs \texttt{infsup\_dense.log}, \texttt{infsup\_lobpcg\_*.log}.
\item \texttt{filter.py MESH N}: method (S), test case A, $\mu=100$; raw vs filtered pressure. Logs \texttt{filter*.log}.
\item \texttt{local\_L.py}: exact checks (Lemma~\ref{lem:K}, Lemma~\ref{lem:bubble}, flux bubble). Log \texttt{local\_L.log}.
\item \texttt{locate\_mode.py MESH N}: localisation of the smallest inf-sup mode.
\end{itemize}

\begin{table}[h]\centering\small
\caption{Inf-sup constants ($k=4$). $\beta_{\rm full}$ on $\Pi_h\cap L^2_0$, $\beta_{\Ql}$ on $\Ql$ (Theorem~\ref{thm:L}), $\beta_{(L)}$ on the (L)-space. Values in parentheses: repeated computations (dense / LOBPCG).}\label{tab:infsup}
\begin{tabular}{lrrrrlrrrr}\toprule
mesh & $N$ & $h_\Gamma$ & $\#\mathcal L$ & $\phi_{\min}$ & $\beta_{\rm full}$ & $\beta_{\rm full}/\phi_{\min}$ & $\beta_{\Ql}$ & $\beta_{(L)}$ & min angle\\\midrule
std & 16 & 0.4595 & 0 & -- & 0.1023 & -- & 0.1023 & 0.1023 & 15.5\\
std & 24 & 0.3204 & 0 & -- & 0.0914 & -- & 0.0914 & 0.0914 & 15.3\\
std & 32 & 0.2444 & 0 & -- & 0.0864 & -- & 0.0864 & 0.0864 & 15.2\\
std & 48 & 0.1650 & 0 & -- & 0.0815 & -- & 0.0815 & 0.0815 & 15.0\\
one & 16 & 0.4595 & 1 & 0.3927 & 0.0583 & 0.149 & 0.1023 & 0.1023 & 15.5\\
one & 24 & 0.3204 & 1 & 0.2318 & 0.0327 & 0.141 & 0.0914 & 0.0914 & 15.3\\
one & 32 & 0.2444 & 1 & 0.1609 & 0.0219 & 0.136 & 0.0864 & 0.0864 & 15.2\\
one & 48 & 0.1650 & 1 & 0.0987 & 0.0128 & 0.130 & 0.0815 & 0.0815 & 15.0\\
alt & 16 & 0.4595 & 8 & 0.3927 & 0.0591 & 0.151 & 0.1023 & 0.1023 & 15.5\\
alt & 24 & 0.3204 & 12 & 0.2318 & 0.0330 (0.0330,0.0330) & 0.142 & 0.0914 & 0.0914 & 15.3\\
alt & 32 & 0.2444 & 16 & 0.1609 & 0.0218 & 0.136 & 0.0864 & 0.0864 & 15.2\\
alt & 48 & 0.1650 & 24 & 0.0987 & 0.0126 & 0.128 & 0.0815 & 0.0815 & 15.0\\
fan & 32 & 0.2444 & 2 & 0.1993 & 0.0316 & 0.159 & 0.0751 & 0.0864 & 10.2\\
fan & 48 & 0.1650 & 2 & 0.1155 & 0.0164 & 0.142 & 0.0815 & 0.0815 & 14.2\\
fan & 64 & 0.1243 & 2 & 0.0801 & 0.0101 & 0.126 & 0.0789 & 0.0789 & 15.0\\
\bottomrule\end{tabular}\end{table}


\begin{table}[h]\centering\small
\caption{Method (S), test case A, $\mu=100$, $k=4$. $\|\pi_h\|$: raw pressure error; $\|\pi_h^f\|$: filtered pressure error (Theorem~\ref{thm:P}(iii)); $L^\infty$ errors on the boundary layer; last column $\|(I-P_{\rm lock})p_h\|$.}\label{tab:filter}
\begin{tabular}{lrrrrrrrrrrrr}\toprule
mesh & $N$ & $h_\Gamma$ & $\#\mathcal L$ & $H^1$ err & rate & $\|\pi_h\|$ & rate & $\|\pi_h^f\|$ & rate & $L^\infty$ & $L^\infty_f$ & perp\\\midrule
one & 16 & 0.4595 & 1 & 5.285e-01 & -- & 5.820 & -- & 2.7590 & -- & 114.7 & 22.64 & 5.125\\
one & 24 & 0.3204 & 1 & 3.068e-01 & 1.51 & 6.462 & -0.29 & 1.4305 & 1.82 & 205.6 & 12.62 & 6.302\\
one & 32 & 0.2444 & 1 & 2.123e-01 & 1.36 & 7.291 & -0.45 & 0.9003 & 1.71 & 308.0 & 12.55 & 7.236\\
one & 48 & 0.1650 & 1 & 1.300e-01 & 1.25 & 8.447 & -0.37 & 0.4863 & 1.57 & 525.0 & 12.38 & 8.433\\
one & 64 & 0.1243 & 1 & 9.376e-02 & 1.16 & 9.160 & -0.29 & 0.3283 & 1.39 & 749.4 & 12.24 & 9.154\\
one & 96 & 0.0831 & 1 & 6.039e-02 & 1.09 & 9.983 & -0.21 & 0.2001 & 1.23 & 1208.4 & 12.07 & 9.981\\
one & 128 & 0.0624 & 1 & 4.464e-02 & 1.06 & 10.447 & -0.16 & 0.1452 & 1.12 & 1673.9 & 11.98 & 10.446\\
alt & 16 & 0.4595 & 8 & 9.030e-01 & -- & 14.020 & -- & 2.2108 & -- & 109.6 & 11.22 & 13.845\\
alt & 24 & 0.3204 & 12 & 6.913e-01 & 0.74 & 17.285 & -0.58 & 1.7052 & 0.72 & 184.7 & 10.54 & 17.200\\
alt & 32 & 0.2444 & 16 & 5.724e-01 & 0.70 & 19.747 & -0.49 & 1.3780 & 0.79 & 264.5 & 9.68 & 19.699\\
alt & 48 & 0.1650 & 24 & 4.443e-01 & 0.64 & 23.429 & -0.44 & 1.0426 & 0.71 & 422.9 & 8.53 & 23.405\\
alt & 64 & 0.1243 & 32 & 3.746e-01 & 0.60 & 26.384 & -0.42 & 0.8683 & 0.65 & 577.1 & 7.93 & 26.370\\
alt & 96 & 0.0831 & 48 & 2.977e-01 & 0.57 & 31.368 & -0.43 & 0.6823 & 0.60 & 877.4 & 7.70 & 31.361\\
alt & 128 & 0.0624 & 64 & 2.543e-01 & 0.55 & 35.653 & -0.45 & 0.5799 & 0.57 & 1173.2 & 7.60 & 35.648\\
fan & 32 & 0.2444 & 2 & 2.148e-01 & -- & 4.132 & -- & 0.8337 & -- & 147.9 & 8.34 & 4.047\\
fan & 48 & 0.1650 & 2 & 1.282e-01 & 1.31 & 4.632 & -0.29 & 0.4127 & 1.79 & 224.9 & 6.37 & 4.614\\
fan & 64 & 0.1243 & 2 & 9.188e-02 & 1.18 & 5.224 & -0.42 & 0.2582 & 1.66 & 325.5 & 6.26 & 5.218\\
fan & 96 & 0.0831 & 2 & 5.935e-02 & 1.09 & 6.123 & -0.39 & 0.1419 & 1.49 & 553.3 & 6.44 & 6.121\\
fan & 128 & 0.0624 & 2 & 4.414e-02 & 1.03 & 6.717 & -0.32 & 0.0975 & 1.31 & 796.1 & 6.62 & 6.717\\
\bottomrule\end{tabular}\end{table}


\paragraph{Observations.}
\begin{itemize}
\item \emph{Theorem~\ref{thm:L} (NUMERICAL confirmation).} On every mesh with locks and $\theta_{\min}\ge14^\circ$,
$\beta_{\Ql}$ and $\beta_{(L)}$ equal $\beta$ of the lock-free mesh \texttt{std} at the same $N$ to 4--5 digits:
the constrained constant does not see the locks at all. This includes mesh (a) (24 locks at $N=48$, alternating) and
\texttt{fan} (two \emph{consecutive} locks sharing an apex). The one exception, \texttt{fan} $N=32$ ($0.0751$), is the
coarse fan with minimum angle $10.2^\circ$; it returns to the \texttt{std} value once $\theta_{\min}\ge14^\circ$.
\item \emph{Theorem~\ref{thm:sharp} (NUMERICAL confirmation).} $\beta_{\rm full}/\phi_{\min}\in[0.126,0.159]$ while
$\phi_{\min}$ falls by a factor $5$; the slow downward drift of the ratio tracks the drift of the lock-free constant
(next item). The smallest unconstrained mode carries $99.1\%$ of its $L^2$ mass on the two triangles of the locked star
(\texttt{locate\_mode.py one 16}). Dense and LOBPCG agree where both were run
(\texttt{one} $N=16,24$: $0.05832/0.03275$ both ways; \texttt{alt} $N=24$). One LOBPCG run (\texttt{alt} $N=24$, (L)-space,
single start, block 4) missed the lowest cluster and returned $0.1125$; the dense value $0.0914$ is tabulated. LOBPCG
values at $N\ge48$ are therefore upper bounds in principle; they agree with the pattern of the dense ones.
\item \emph{The lock-free constant drifts:} $0.102,0.091,0.086,0.081,0.079$ for $N=16,24,32,48,64$. The differences shrink
roughly linearly in $\hG$ (extrapolated limit $\approx0.07$). The minimising modes are 4-fold degenerate and sit in the
\emph{outer} layer next to $\partial R$, not at $\Gamma_h$ (\texttt{locate\_mode.py std 16}). This is consistent with the
paper's uniform \texttt{lem:infsup}, but a referee could ask about it. It is independent of the locking question.
\item \emph{Pressure.} Raw error: \texttt{one} $5.8\to10.4$ (bounded, $\to\approx12$ by an $a-b\hG$ fit); mesh (a)
$14\to35.7$, local rate $-0.45\to-1/2$; \texttt{fan} $4.1\to6.7$. $L^\infty$ on the layer $\propto\hG^{-1}$
($\hG\,\norm{\pi_h}_{L^\infty}\approx93,100,104.5$ for \texttt{one} and $71.7,72.9,73.2$ for mesh (a), $N=64,96,128$). Filtered error rates:
$1.12$ (\texttt{one}), $0.57$ (mesh (a)), $1.31$ (\texttt{fan}), close to the velocity rates $1.06$, $0.55$, $1.04$.
The $L^\infty$ error of the filtered pressure on the layer stays $O(1)$, as it must: an $L^2$ error $\asymp\hG|\dn u|$ on a
patch of area $\hG^2$.
\end{itemize}

\section{Suggested use in the paper (not made)}
Replace ``S4 (PROVED MODULO (L))'' by Theorem~\ref{thm:S4} + Corollary~\ref{cor:two}; state (H) as ``(M0), (M1), and (M2) at all
vertices except $\Gamma_h$ vertices with two triangles''. Add Theorem~\ref{thm:P}: with locks the \emph{velocity} is
covered, the \emph{raw pressure} is not (non-convergent), and the one-line filter $P_{\rm lock}$ restores it. The
key new device is the exact identity $D_z(w-z)=(1,1)$ (Lemma~\ref{lem:K}(d)): the radial field
$(w-z)\lambda_w\lambda_z^2$ produces equal vertex divergences at a locked vertex at $\phi$-independent cost.
Literature caveat: Guzm\'an--Scott and Scott--Vogelius treat singular and non-singular vertices with constants
$\propto\Theta^{-1}$; for nearly singular vertices, pressure non-robustness with robust velocities has been studied in
the elasticity/Stokes literature (e.g.\ work of Ainsworth and Parker on nearly singular vertices); I did not verify
whether a statement equivalent to Theorem~\ref{thm:sharp} appears there, so priority should be checked before claiming it.

\end{document}
